% Part: turing-machines
% Chapter: machines-computations
% Section: well-behaved-machines

\documentclass[../../../include/open-logic-section]{subfiles}

\begin{document}

\olfileid{tur}{mac}{dis}
\olsection{منضبط ماشينونه}

\begin{explain}
په برخه~\olref[hal]{sec} کښې مو هغه ټيورينګ ماشينونه وڅېړل چې يو
ټاکلی درېدنی حالت~$h$ لري---داسې ماشينونه، که بېخي درېږي، هرومرو په
حالت~$h$ کښې درېږي۔ په دې معنا د يوه درېدني حالت لرونکي ماشينونه د
هغو ټيورينګ ماشينونو په پرتله ډېر ``منضبط'' دي چې موږ ئې په عمومي
ډول منو۔ نور بنديزونه هم د ټيورينګ ماشينونو پر چلند لګولے شو۔ مثلاً،
موږ ټيورينګ ماشينونه له دې هم نۀ دي منع کړي چې کله د پټې د پای نښه
پر مربع~$0$ پاکه کړي، يا له مربع~$0$ څخه کيڼ لوري ته د تلو هڅه وکړي۔
(زموږ تعريف وايي چې په دې صورت کښې سر بس پر مربع~$0$ پاتې کېږي؛ په
نورو تعريفونو کښې ماشين درېږي۔) همداراز کله ناکله دا فرض هم ګټور وي
چې ټيورينګ ماشين، که بېخي درېږي، پر مربع~$1$ ودرېږي۔
\end{explain}

\begin{defn}\ollabel{defn:disciplined}
ټيورينګ ماشين~$M$ هله او يوازې هله \emph{منضبط} دے چې
\begin{enumerate}
    \item يو ټاکلی يوازينی درېدنی حالت~$h$ ولري،
    \item که بېخي درېږي، د مربع~$1$ د لوستلو پر وخت ودرېږي،
    \item پر مربع~$0$ د~$\TMendtape$ نښه هېڅکله پاکه نۀ کړي، او
    \item له مربع~$0$ څخه کيڼ لوري ته د تلو هڅه هېڅکله ونۀ کړي۔
\end{enumerate}
\end{defn}

\begin{explain}
موږ مخکښې بحث کړے چې هر ټيورينګ ماشين په داسې ماشين بدلېدے شي چې
هماغه چلند او يو ټاکلی درېدنی حالت ولري۔ دا کار يوازې د نوي حالت~$h$
په زياتولو، او د هر داسې جوړې~$\tuple{q,\sigma}$ دپاره د لارښوونې
$\delta(q, \sigma) = \tuple{h, \sigma, N}$ په زياتولو کېږي چې اصلي
$\delta$ پرې تعريف شوې نۀ وي۔ که څه هم ثبوت ئې ستړے کوونکے دے، خو
دا رښتيا ده چې هر ټيورينګ ماشين~$M$ په داسې منضبط ټيورينګ ماشين~$M'$
بدلېدے شي چې پر هماغو ننوتونو درېږي او هماغه وت پيدا کوي۔ مثلاً، که
ټيورينګ ماشين ودرېږي او پر مربع~$1$ نۀ وي، داسې لارښوونې ور زياتولے
شو چې سر تر هغه کيڼ لوري ته بوځي څو د پټې د پای نښه پيدا کړي، بيا يو
مربع ښي لوري ته لاړ شي، او بيا ودرېږي۔ د نورو شرطونو د پوره کولو پر
لار به تاسو پخپله فکر وکړئ۔
\end{explain}

\begin{ex}
    په~\olref{fig:adder-disc} کښې د جمع ماشين چې په
    \olref[una]{ex:adder} کښې و، په منضبط ماشين بدلوو۔
    \begin{figure}\[
\begin{tikzpicture}[->,>=stealth',shorten >=1pt,auto,node distance=2.8cm,
                    semithick]
  \tikzstyle{every state}=[fill=none,draw=black,text=black]

  \node[initial,state]         (A)              {$q_0$};
  \node[state]         (B) [right of=A] {$q_1$};
  \node[state]         (C) [below right of=B] {$q_2$};
  \node[state]         (D) [below left of=C] {$q_3$};
  \node[state]         (H) [left of=D] {$h$};

  \path (A) edge node {\TMtrans{\TMblank}{\TMstroke}{\TMstay}} (B)
           edge [loop below] node {\TMtrans{\TMstroke}{\TMstroke}{\TMright}} (B)
        (B) edge [loop below] node {\TMtrans{\TMstroke}{\TMstroke}{\TMright}} (B)
            edge node {\TMtrans{\TMblank}{\TMblank}{\TMleft}} (C)
        (C) edge node {\TMtrans{\TMstroke}{\TMblank}{\TMleft}} (D)
        (D) edge [loop above] node {\TMtrans{\TMstroke}{\TMstroke}{\TMleft}} (D)
        edge node {\TMtrans{\TMendtape}{\TMendtape}{\TMright}} (H);
\end{tikzpicture}
\]\caption{يو منضبط جمع کوونکے ماشين}
\ollabel{fig:adder-disc}
\end{figure}
\end{ex}

\begin{prop}\ollabel{prop:disciplined} د هر ټيورينګ ماشين~$M$ دپاره
داسې منضبط ټيورينګ ماشين~$M'$ شته چې که~$M$ د وت~$O$ سره درېږي،
هغه هم د وت~$O$ سره درېږي؛ او که~$M$ نۀ درېږي، هغه هم نۀ درېږي۔ په
ځانګړې توګه، هره تابع $f\colon\Nat^n \to \Nat$ چې د ټيورينګ ماشين
په وسيله محاسبه کېدونکې وي، د منضبط ټيورينګ ماشين په وسيله هم محاسبه
کېدونکې ده۔
\end{prop}

\begin{prob}\label{tur:mac:dis:prob:disc-succ}
داسې منضبط ماشين ورکړئ چې $f(x) = x+1$ محاسبه کړي۔
\end{prob}


\begin{prob}\label{tur:mac:dis:prob:copier}
داسې منضبط ماشين پيدا کړئ چې پر ننوت $\TMstroke^n$ پيل شي او وت
$\TMstroke^n \concat \TMblank \concat \TMstroke^n$ پيدا کړي۔
\end{prob}

\end{document}
